\chapter{Formal Framework}\label{chap:framework}

We consider the problem of optimizing the makespan of partially ordered plans. We first define a classical planning problem and plan in STRIPS notation, and then extend it to POCL problems and plans. We follow the notation used by~\cite{Bercher2020POPvsPOCL}.\\


A classical planning problem $\mathcal{P}$ is given by a tuple $(V,A,s_I,g)$, where $V$ is a finite set of propositional state variables, $A$ is a finite set of \textit{actions}, $s_I\in 2^V$ is the initial state, and $g\subseteq V$ is the goal description. An \textit{action} $a\in A$ is a tuple $(pre, add, del)$, referred to as $pre(a)$, $add(a)$, and $del(a)$, where $pre, add, del \subseteq V$. These sets represent an action's \textit{precondition}, \textit{add}, and \textit{delete} lists. As is standard, an action $a\in A$ is executable in a state $s\in 2^V$ if and only if $pre(a)\subseteq s$. If $a$ is executable in $s$, the state transition function $\gamma : A \times 2^F \rightarrow 2^F$ returns its successor state after executing $a$: $\gamma (a,s)=(s\setminus del(a))\cup add(a)$. A sequence of actions $\overline{a}=(a_0,a_1,\ldots,a_n)$ is executable in a state $s_0$ if there exists a state sequence $(s_1,s_2,\ldots ,s_{n+1})$ such that $\gamma (a_{i-1},s{i-1})=s_i$ holds for all $1\leq i \leq n+1$. The state $s_{n+1}$ is called the \textit{result} of executing the sequence. If a state $s\in 2^V$ follows $g\subseteq s$ then it is called a \textit{goal state}.\\


Classical solutions to planning problems are total-order plans, as defined below.\newline

\textbf{Definition 1.} An action sequence $\overline{a}$ is called a \textit{total-order plan} $\overline{P}$ (also \textit{total-order solution}) for a planning problem if and only if it is executable in $s_I$ and results in a goal state.\\


We will next extend this definition to POCL plans. POCL planning does not produce total-order plans, and instead relies on \textit{orderings} and \textit{causal links} to produce \textit{partially ordered plans}. When using POCL search algorithms to solve planning problems, search nodes are \textit{partial plans}, as opposed to classical planning where each search node is a \textit{state}. We will also refer to these partial plans as \textit{plans} for short when it is clear from context that we are referring to a partial plan and not a solution.\\


A partial plan $P$ is given by a tuple $(PS,\prec , CL)$, where $PS$ is a finite set of \textit{plan steps} $ps=(l,a)$ with $l$ being a unique label and $a\in A$ being an action, $\prec$ is a \textit{partial order} on $PS$, and $CL$ is a finite set of \textit{causal links}. Plan steps require unique labels to differentiate between multiple occurrences of the same action within a partial plan. As with actions, we use $pre(ps),add(ps),$ and $del(ps)$ to refer to the precondition, add, and delete lists of a plan step's action. A partial order is a finite set of orderings $ps_n \prec ps_m$ between steps. A causal link $cl=(ps_p, v, ps_c)\in CL$ where $ps_p \in PS, v\in V, ps_c \in PS$, is used to formally represent a causal dependency between plan steps. This causal link indicates that the precondition $v$ of the \textit{consumer} plan step $ps_c$ is \textit{supported} by the \textit{producer} plan step $ps_p$, meaning that $v\in add(ps_p)$ is currently ensuring that the condition $v\in pre(ps_c)$ is met. A situation where a plan step $ps_t$ with $v\in del(ps_t)$ may be ordered between $ps_p$ and $ps_c$ is called a \textit{causal threat}, and $ps_t$ is called a \textit{threatening} plan step. The variable $v$ is said to be \textit{protected} by its causal link as any causal threat to this protection is raised as a \textit{flaw} in the plan. We always assume for simplicity and without loss of generality that $ps_p\prec ps_c$.\\


In POCL planning, the problem's initial state and goal description are encoded using two artificial actions, labelled \textit{init} and \textit{goal}, where $add(init)=\{v\in s_I: v \text{ is true}\}$, and $pre(goal)=g$. The sets $pre(init),del(init),add(goal),$ and $del(goal)$ are all empty. In all partial plans, \textit{init} is ordered before all other steps and \textit{goal} is ordered after them. \\


We can now define POCL solutions to planning problems.\newline

\textbf{Definition 2.} A partial POCL plan $P$ is called a \textit{POCL plan} (also \textit{POCL solution}) for a planning problem if and only if every precondition is supported by a causal link and there are no causal threats.\\


It is common knowledge that this criteria for a POCL plan is sufficient to imply that every linearization of its steps is executable, and hence a POCL plan can represent up to an exponential set of total-order solutions. Note that the inverse is not true, i.e every linearization being executable is \textit{not} a sufficient condition for a partial POCL plan to be a POCL plan, as proved by~\cite{Bercher2020POPvsPOCL}.